\documentclass[10pt,a4paper]{amsart}
\usepackage[margin=19mm]{geometry}
\usepackage{amsmath,amssymb,mathtools}
\usepackage[scr=rsfs,cal=euler]{mathalfa}
\usepackage{xfrac}
\usepackage[unicode,hidelinks,bookmarks=false]{hyperref}
\newcommand{\ie}{i.e.\ }
\newcommand{\cf}{cf.\ }
\newcommand{\resp}{respectively\ }
\newcommand{\Subsection}{\subsection}
\providecommand{\nobreakdash}{\nobreak\hbox{-}}
\setlength{\emergencystretch}{2em}
\multlinegap0pt
\usepackage{xcolor}
\newtheorem{lemma}{Lemma}
\newtheorem{theorem}{Theorem}
\newcommand{\blambda}{{\boldsymbol{\lambda}}}
\newcommand{\ee}{{\mathbb E}}
\newcommand{\odim}{\overline{\dim}}
\newcommand{\udim}{\underline{\dim}}
\title{The Attainable almost sure LARGE dimensions}
\author{Kathryn E. Hare, Franklin Mendivil}
\date{Source: arXiv:2604.09899v1 (CC BY 4.0)}
\begin{document}
\maketitle

\noindent\emph{Source and licence.} The Attainable almost sure LARGE dimensions. Version arXiv:2604.09899v1 (CC BY 4.0), distributed by arXiv under the Creative Commons Attribution 4.0 International licence, http://creativecommons.org/licenses/by/4.0/, as recorded in the arXiv licence metadata for this exact version and shown on the abstract page https://arxiv.org/abs/2604.09899v1. Full licence notice: \texttt{licenses/arxiv-2604.09899-CC-BY-4.0.txt}.\par\smallskip

\noindent\textit{Editorial scope.} Counted unit: the article's theorem thm:G(t)=t (source0.tex lines 686--696). The article's complete original proof of it is delivered with it (source0.tex lines 698--803, one \texttt{proof} environment of 106 lines). No mathematics is written, repaired or re-derived: the statement and the proof are the article's own lines, in the article's own notation, and no retained line is substituted or re-flowed.  Retained uncounted as the source-local context the proof needs: lines 627--637.  Nothing was omitted from any retained span, so the package carries no omission record.  No retained line contains a citation, so the wrapper carries no bibliography and no bibliography entry.  No retained line contains a figure environment, an image-inclusion command or drawing code, so no image is delivered and the package declares no asset dependency.  Editorial formatting only, in addition: an AMS \texttt{amsart} preamble that restates the article's own declarations and macros used by the retained lines, namely the environments \texttt{lemma}, \texttt{theorem} on separate counters, together with the standard packages those lines need.  The retained environments print in this document as Lemma 1, Theorem 1 in order of appearance, whereas the article prints the same items with the numbering of its own full document.  The complete native source of the article as distributed in the arXiv e-print for this exact version is bundled unchanged as \texttt{sources/198233/source0.tex}.
\begin{lemma}
\label{lem:H-function}
Let $\chi _{n}:\Lambda \to \{1,2, \ldots, K\}$ be iid rv's  and let 
\begin{equation*}
    H(\chi )=\frac{\ee(\log p(\chi ,\lambda ))}{\ee(\log a(\chi,\lambda ))}.
\end{equation*}
Then 
\begin{equation*}
    \inf_{\theta }G^{\prime }(\theta )\leq H(\chi )\leq \sup_{\theta }G(\theta ).
\end{equation*}
\end{lemma}

\begin{theorem} \label{thm:G(t)=t}

(i) $\sup_{\chi }H(\chi )=\sup_{\theta }G(\theta )=\odim_\Phi\mu_\blambda $ a.s.

(ii) $t= \sup_{\chi }H(\chi )$ if and only if $G(t)=t$

(iii) $\inf_\chi H(\chi) = \inf_\theta G'(\theta) = \udim_\Phi \mu_\blambda$ a.s.

(iv) $t = \inf_\chi H(\chi)$ if and only if $G'(t) = t$
 
\end{theorem}

\begin{proof}
Statements (i) and (iii) are immediate from Lemma \ref{lem:H-function}.

(ii) First, suppose $t=G(t)$. 
Then $t\leq \sup_{\theta }G(\theta)=\sup_{\chi }H(\chi )$ so we only need prove the other inequality.

To shorten notation, we will write $m_{n}=m(t,\lambda _{n})$ and $\chi _{n}=\chi(\lambda _{n})$.

We first observe that for any $\chi$  
\begin{equation*}
      \lim_{N}\log \prod_{n=1}^{N}a(\chi _{n})^{t/N}=
    t\lim_{N}\frac{1}{N}\sum_{n=1}^{N}\log a(\chi _{n})=
   t\ee(\log a(\chi ))\text{ for a.a. } \lambda 
\end{equation*}
and since the exponential function is continuous it follows that $\lim_{N}\prod_{n=1}^{N}a(\chi _{n})^{t/N}$ exists a.s.  
Furthermore, this limit is bounded away from $0$ as the contraction factors are bounded away from $0$. 
Similarly, $\lim_{N}\prod_{n=1}^{N}p(\chi _{n})^{1/N}$ exists a.s.

Since $t=G(t)$ we have $t\ee(\log a(m))=\ee(\log p(m))$.
Consequently, almost surely
\begin{equation*}
    \lim_{N}\log \prod_{n=1}^{N}a(m_{n})^{t/N}
  = \lim_{N}\frac{t}{N}\sum_{n=1}^{N}\log a(m_{n})
  =\lim_{N}\frac{1}{N}\sum_{n=1}^{N}\log p(m_{n})
  =\lim_{N}\log \prod_{n=1}^{N}p(m_{n})^{1/N}
\end{equation*}
and therefore 
\begin{equation*}
     \lim_{N}\prod_{n=1}^{N}a(m_{n})^{t/N}=\lim_{N}\prod_{n=1}^{N}p(m_{n})^{1/N},
\end{equation*}
equivalently,  
\begin{equation*}
      1=\lim_{N}\left( \prod_{n=1}^{N}\frac{p(m_{n})}{a(m_{n})^{t}}\right) ^{1/N}.
\end{equation*}
But by the definition of $m$, 
\begin{equation*}
     \frac{p(m_{n})}{a(m_{n})^{t}}\leq \frac{p(\chi _{n})}{a(\chi _{n})^{t}}\text{for all }\chi ,
\end{equation*}
thus
\begin{equation*}
   1\leq \lim_{N}\left( \prod_{n=1}^{N}\frac{p(\chi _{n})}{a(\chi _{n})^{t}}\right) ^{1/N}
\end{equation*}
and that implies
\begin{equation*}
   \lim_{N}\prod_{n=1}^{N}a(\chi _{n})^{t/N}\leq 
           \lim_{N}\prod_{n=1}^{N}p(\chi_{n})^{1/N}.
\end{equation*}
Taking logarithms we deduce that
\begin{equation*}
   t \ee(\log a(\chi ))=\lim_{N}\frac{t}{N}\sum_{n=1}^{N}\log a(\chi_{n})
     \leq \lim_{N}\frac{1}{N}\sum_{n=1}^{N}\log p(\chi _{n})=\ee(\log p(\chi ))\text{ a.s.}
\end{equation*}
hence  
\begin{equation*}
    t\geq \ee(\log p(\chi ))/\ee(\log a(\chi ))=H(\chi ).
\end{equation*}
Since this holds for all $\chi $ it follows that $t\geq \sup_{\chi }H(\chi )$
as we desired to show.

\medskip 

Now assume that $t=\sup_{\chi }H(\chi )=\sup_{\theta }G(\theta )$, but that $t\neq G(t)$. 
In that case, $G(t)=t-\varepsilon $ for some $\varepsilon >0$.
But then for a.a. $\lambda $ and for $N$ sufficiently large 
\begin{equation*}
   t-\varepsilon \geq 
    \frac{\frac{1}{N}\sum_{n=1}^{N}\log p(m_{n})}{\frac{1}{N}\sum_{n=1}^{N}\log a(m_{n})}
     -\frac{\varepsilon }{2}.
\end{equation*}
Hence $(t-\varepsilon /2)\sum_{n=1}^{N}\log a(m_{n})\leq \sum_{n=1}^{N}\log p(m_{n})$ 
for large enough $N,$ whence $\prod_{n=1}^{N}a(m_{n})^{t-\varepsilon /2}\leq \prod_{n=1}^{N}p(m_{n})$, 
or, equivalently,
\begin{equation*}
    \frac{\prod_{n=1}^{N}a(m_{n})^{t}}{\prod_{n=1}^{N}p(m_{n})}
        \leq
        \prod_{n=1}^{N}a(m_{n})^{\varepsilon /2}.
\end{equation*}
It follows from the definition of $m$ that for any $\chi $ we have 
\begin{equation*}
         \prod_{n=1}^{N}a(\chi _{n})^{t}\leq \prod_{n=1}^{N}a(m_{n})^{\varepsilon/2}\prod_{n=1}^{N}p(\chi _{n})
\end{equation*}
and thus
\begin{equation*}
     \frac{t}{N}\sum_{n=1}^{N}\log a(\chi _{n})\leq 
         \frac{\varepsilon /2}{N}\sum_{n=1}^{N}\log a(m_{n})+
        \frac{1}{N}\sum_{n=1}^{N}\log p(\chi _{n}).
\end{equation*}
Taking limits we deduce that 
\begin{equation*}
     t\ee(\log a(\chi ))\leq (\varepsilon /2)\ee(\log a(m))+\ee(\log p(\chi )).
\end{equation*}

Since the contraction factors are bounded away from $0,1$ there is some constant $C>0$ 
such that $\ee(\log a(m))/\ee(\log a(\chi))\geq C > 0$ for all $\chi$.
Thus 
\begin{equation*}
     t\geq 
       \frac{\varepsilon }{2}\frac{\ee(\log a(m))}{\ee(\log a(\chi ))}+
       \frac{\ee(\log p(\chi ))}{\ee(\log a(\chi ))}
    \geq   \frac{C\varepsilon }{2}+H(\chi )\text{ for all }\chi 
\end{equation*}%
and therefore $t\geq C\varepsilon /2+\sup_{\chi }H(\chi )$ which is a
contradiction.

The proof of  (iv) is similar using the obvious modifications.
\end{proof}

\end{document}
